% answers.tex —— 1977 年·北京卷·理科 · 参考答案与详细解析片段
% 由 tools/generate_projects.py 依来源总账自 1977/1977_数学真题_北京理_试题.md + 1977/1977_数学真题_北京理_解析正文.md + 1977/1977_数学真题_北京理_答案要点.md 生成，勿手改（改源 MD 后重跑）


\begin{gktip}
固定源为 A3 单页题面，未附答案页；下列推导由编者独立完成，非官方标准答案。第 11、12 题保留为原卷附加题。
\end{gktip}

\begin{answer}
\gksolhead{1}{}

\textbf{题目}：

\Needspace{5\baselineskip}
解方程：

\nopagebreak[4]
\[
\sqrt{x-1}=3-x.
\]

\Needspace{4\baselineskip}
\textbf{解析}

\Needspace{5\baselineskip}
由根式有意义及等式右端非负，须有 $x\geqslant1$ 且 $x\leqslant3$。平方得

\nopagebreak[4]
\[
x-1=(3-x)^2
\Longleftrightarrow x^2-7x+10=0
\Longleftrightarrow (x-2)(x-5)=0.
\]

在定义域内只保留 $x=2$。

\gkanswer{答案：$x=2$。}

\end{answer}
\begin{answer}
\gksolhead{2}{}

\textbf{题目}：

\Needspace{5\baselineskip}
计算：

\nopagebreak[4]
\[
2^{-\frac12}+\frac{2^0}{\sqrt2}+\frac1{\sqrt2-1}.
\]

\Needspace{4\baselineskip}
\textbf{解析}

\nopagebreak[4]
\[
2^{-\frac12}=\frac1{\sqrt2},
\qquad
\frac{2^0}{\sqrt2}=\frac1{\sqrt2},
\qquad
\frac1{\sqrt2-1}=\sqrt2+1.
\]

所以原式为 $\sqrt2+\sqrt2+1=1+2\sqrt2$。

\gkanswer{答案：$1+2\sqrt2$。}

\end{answer}
\begin{answer}
\gksolhead{3}{}

\textbf{题目}：

已知 $\lg2=0.3010$，$\lg3=0.4771$，求 $\lg\sqrt{45}$。

\Needspace{4\baselineskip}
\textbf{解析}

\nopagebreak[4]
\[
\lg\sqrt{45}=\frac12\lg45
=\frac12\left(\lg10-\lg2+2\lg3\right)
=\frac12(1-0.3010+0.9542)=0.8266.
\]

固定源图像中的根号横线覆盖“45”，故按题面采用 $0.8266$。

\gkanswer{答案：$\lg\sqrt{45}=0.8266$。}

\end{answer}
\begin{answer}
\gksolhead{4}{}

\textbf{题目}：

\Needspace{5\baselineskip}
证明：

\nopagebreak[4]
\[
(1+\tan\alpha)^2=\frac{1+\sin2\alpha}{\cos^2\alpha}.
\]

\Needspace{4\baselineskip}
\textbf{解析}

\nopagebreak[4]
\[
1+\sin2\alpha
=1+2\sin\alpha\cos\alpha
=\sin^2\alpha+\cos^2\alpha+2\sin\alpha\cos\alpha
=(\sin\alpha+\cos\alpha)^2.
\]

\Needspace{5\baselineskip}
另一方面，

\nopagebreak[4]
\[
(1+\tan\alpha)^2
=\left(\frac{\cos\alpha+\sin\alpha}{\cos\alpha}\right)^2
=\frac{(\sin\alpha+\cos\alpha)^2}{\cos^2\alpha}.
\]

两式相等，得所证结论（在 $\cos\alpha\ne0$ 时原式有意义）。

\gkanswer{答案：见解析版}

\end{answer}
\begin{answer}
\gksolhead{5}{}

\textbf{题目}：

求过两直线 $x+y-7=0$ 和 $3x-y-1=0$ 的交点且过 $(1,1)$ 点的直线方程。

\Needspace{4\baselineskip}
\textbf{解析}

\Needspace{5\baselineskip}
两直线交点满足

\nopagebreak[4]
\[
x+y=7,\qquad3x-y=1,
\]

\Needspace{5\baselineskip}
相加得 $4x=8$，所以交点为 $(2,5)$。过 $(2,5)$ 与 $(1,1)$ 的直线斜率为 $4$，故

\nopagebreak[4]
\[
y-1=4(x-1)\Longleftrightarrow4x-y-3=0.
\]

\gkanswer{答案：$4x-y-3=0$。}

\end{answer}
\begin{answer}
\gksolhead{6}{}

\textbf{题目}：

某工厂今年七月份的产值为 $100$ 万元，以后每月产值比上月增加 $20\%$，问今年七月份到十月份总产值是多少？

\Needspace{4\baselineskip}
\textbf{解析}

\Needspace{5\baselineskip}
七月至十月共四个月，产值构成首项 $100$、公比 $1.2$ 的等比数列。因此总产值为

\nopagebreak[4]
\[
100(1+1.2+1.2^2+1.2^3)=436.8\text{（\text{万元}）}.
\]

\gkanswer{答案：$436.8$ 万元。}

\end{answer}
\begin{answer}
\gksolhead{7}{}

\textbf{题目}：

已知二次函数 $y=x^2-6x+5$。

（1） 求出它的图象的顶点坐标和对称轴方程；

（2） 画出它的图象；

（3） 分别求出它的图象和 $x$ 轴、$y$ 轴的交点坐标。

\Needspace{4\baselineskip}
\textbf{解析}

\gksubpart{（1）}

\nopagebreak[4]
\[
y=x^2-6x+5=(x-3)^2-4,
\]

所以顶点为 $(3,-4)$，对称轴为 $x=3$。

\gksubpart{（2）}

抛物线开口向上，过 $(1,0)$、$(5,0)$、$(0,5)$，顶点为 $(3,-4)$；按这些点及对称轴即可作图。
\nopagebreak[4]
\begin{center}% 1977 年北京卷理科第 7 题：二次函数图象
\begin{tikzpicture}[scale=0.62,>=stealth]
  \draw[->] (-1.1,0)--(7.1,0) node[right] {$x$};
  \draw[->] (0,-5.2)--(0,5.0) node[above] {$y$};
  \draw[thick,domain=0.2:5.8,smooth,variable=\x]
    plot ({\x},{(\x-3)^2-4});
  \fill (3,-4) circle (1.3pt) node[below right] {$\left(3,-4\right)$};
  \fill (1,0) circle (1.2pt) node[above left] {$1$};
  \fill (5,0) circle (1.2pt) node[above right] {$5$};
  \fill (0,5) circle (1.2pt) node[left] {$5$};
  \draw[dashed] (3,-4.6)--(3,4.5) node[above] {$x=3$};
\end{tikzpicture}
\end{center}
\nopagebreak[4]
\gksubpart{（3）}

\Needspace{5\baselineskip}
令 $y=0$，得

\nopagebreak[4]
\[
(x-3)^2-4=0\Longleftrightarrow x=1\text{ \text{或} }5,
\]

故与 $x$ 轴交于 $(1,0)$、$(5,0)$。令 $x=0$，得 $y=5$，故与 $y$ 轴交于 $(0,5)$。

\gkanswer{答案（1）：顶点为 $(3,-4)$，对称轴为 $x=3$。}

\gkanswer{答案（2）：见解析版。}

\gkanswer{答案（3）：与 $x$ 轴交于 $(1,0)$、$(5,0)$，与 $y$ 轴交于 $(0,5)$。}

\end{answer}
\begin{answer}
\gksolhead{8}{}

\textbf{题目}：

一只船以 $20$ 海里/小时的速度向正东航行，起初船在 $A$ 处看见一灯塔 $B$ 在船的北 $45^\circ$ 东方向，一小时后船在 $C$ 处看见这个灯塔在船的北 $15^\circ$ 东方向，求这时船和灯塔的距离 $CB$。

\Needspace{4\baselineskip}
\textbf{解析}

船一小时向正东航行 $20$ 海里，所以 $AC=20$。在三角形 $ABC$ 中，$\angle A=45^\circ$；从 $C$ 看，$CA$ 指向正西，而 $CB$ 为北偏东 $15^\circ$，故 $\angle C=105^\circ$，从而 $\angle B=30^\circ$。

\Needspace{5\baselineskip}
由正弦定理

\nopagebreak[4]
\[
\frac{CB}{\sin45^\circ}=\frac{AC}{\sin30^\circ},
\]

\Needspace{5\baselineskip}
得到

\nopagebreak[4]
\[
CB=20\cdot\frac{\sin45^\circ}{\sin30^\circ}=20\sqrt2\text{（\text{海里}）}.
\]

\gkanswer{答案：$CB=20\sqrt2$ 海里。}

\end{answer}
\begin{answer}
\gksolhead{9}{}

\textbf{题目}：

\Needspace{5\baselineskip}
一个圆内接三角形 $ABC$，$\angle A$ 的平分线交 $BC$ 于 $D$，交外接圆于 $E$，求证：

\nopagebreak[4]
\[
AD\cdot AE=AC\cdot AB.
\]

\Needspace{4\baselineskip}
\textbf{解析}

\Needspace{5\baselineskip}
由 $AD$ 是 $\angle A$ 的平分线，角平分线定理给出

\nopagebreak[4]
\[
\frac{DB}{DC}=\frac{AB}{AC}.
\]

\Needspace{5\baselineskip}
又因 $D$ 是弦 $BC$ 与弦 $AE$ 的交点，交点弦定理给出

\nopagebreak[4]
\[
DA\cdot DE=DB\cdot DC.
\]

\Needspace{5\baselineskip}
设 $AD=u$、$DE=v$，则 $AE=u+v$。由上式及角平分线定理的等价形式，可得

\nopagebreak[4]
\[
AD\cdot AE=AD^2+AD\cdot DE
=AD^2+DB\cdot DC=AB\cdot AC,
\]

\Needspace{5\baselineskip}
其中最后一步是角平分线长度关系（也可由 Stewart 定理直接验证）。故

\nopagebreak[4]
\[
AD\cdot AE=AC\cdot AB.
\]

\gkanswer{答案：见解析版}

\end{answer}
\begin{answer}
\gksolhead{10}{}

\textbf{题目}：

\Needspace{5\baselineskip}
当 $m$ 取哪些值时，直线 $y=x+m$ 与椭圆

\nopagebreak[4]
\[
\frac{x^2}{16}+\frac{y^2}{9}=1
\]

有一个交点？有两个交点？没有交点？当它们有一个交点时，画出它的图象。

\begin{gktip}
\Needspace{0.25\textheight}
\textbf{附加题}
\end{gktip}

\Needspace{5\baselineskip}
\textbf{附加题 11．}

\nopagebreak[4]
\[
f(x)=
\begin{cases}
x^2\sin\dfrac\pi x,&x\ne0,\\
0,&x=0,
\end{cases}
\]

求函数 $f(x)$ 的导数。

\textbf{附加题 12．}

\begin{itemize}
\item \textbf{（1）} 试用 $\varepsilon-\delta$ 语言叙述“函数 $f(x)$ 在点 $x=x_0$ 处连续”的定义。
\end{itemize}

\begin{itemize}
\item \textbf{（2）} 试证明：若 $f(x)$ 在点 $x=x_0$ 处连续，且 $f(x_0)>0$，则存在一个 $x_0$ 的邻域 $(x_0-\delta,x_0+\delta)$，在这个邻域内，处处有 $f(x)>0$。
\end{itemize}

\Needspace{4\baselineskip}
\textbf{解析}

\Needspace{5\baselineskip}
把直线 $y=x+m$ 代入椭圆，得

\nopagebreak[4]
\[
\frac{x^2}{16}+\frac{(x+m)^2}{9}=1,
\]

\Needspace{5\baselineskip}
即

\nopagebreak[4]
\[
25x^2+32mx+16m^2-144=0.
\]

\Needspace{5\baselineskip}
判别式为

\nopagebreak[4]
\[
\Delta=(32m)^2-4\cdot25(16m^2-144)=576(25-m^2).
\]

所以 $|m|<5$ 时有两个交点，$|m|=5$ 时有一个交点，$|m|>5$ 时没有交点。相切时可取 $m=5$ 作图（另一条相切线为 $m=-5$）。
\nopagebreak[4]
\begin{center}% 1977 年北京卷理科第 10 题：椭圆与相切直线示意
\begin{tikzpicture}[scale=0.65,>=stealth]
  \draw[->] (-5.2,0)--(5.2,0) node[right] {$x$};
  \draw[->] (0,-4.2)--(0,4.2) node[above] {$y$};
  \draw[thick,domain=0:360,smooth,variable=\t]
    plot ({4*cos(\t)},{3*sin(\t)});
  \draw[thick] (-5,0)--(-1.0,4.0) node[above right] {$y=x+5$};
  \fill (-1,4) circle (1.3pt) node[above right] {切点};
  \node[below right] at (3.7,2.7) {$\dfrac{x^2}{16}+\dfrac{y^2}{9}=1$};
\end{tikzpicture}
\end{center}
\nopagebreak[4]
\gksubpart{附加题 11}

\Needspace{4\baselineskip}
\textbf{解析}

\Needspace{5\baselineskip}
当 $x\ne0$ 时，

\nopagebreak[4]
\[
f'(x)=2x\sin\frac\pi x+x^2\cos\frac\pi x\left(-\frac\pi{x^2}\right)
=2x\sin\frac\pi x-\pi\cos\frac\pi x.
\]

\Needspace{5\baselineskip}
在 $x=0$ 处由导数定义，

\nopagebreak[4]
\[
f'(0)=\lim_{h\to0}\frac{h^2\sin(\pi/h)}h
=\lim_{h\to0}h\sin(\pi/h)=0.
\]

\Needspace{5\baselineskip}
故

\nopagebreak[4]
\[
f'(x)=
\begin{cases}
2x\sin\dfrac\pi x-\pi\cos\dfrac\pi x,&x\ne0,\\
0,&x=0.
\end{cases}
\]

\gksubpart{附加题 12}

\Needspace{4\baselineskip}
\textbf{解析}

\Needspace{5\baselineskip}
（1）函数 $f$ 在 $x=x_0$ 处连续，是指：对任意 $\varepsilon>0$，都存在 $\delta>0$，使得当 $|x-x_0|<\delta$ 时，有

\nopagebreak[4]
\[
|f(x)-f(x_0)|<\varepsilon.
\]

\Needspace{5\baselineskip}
（2）取 $\varepsilon=f(x_0)/2>0$。由连续性，存在 $\delta>0$，使 $|x-x_0|<\delta$ 时

\nopagebreak[4]
\[
|f(x)-f(x_0)|<\frac{f(x_0)}2.
\]

\Needspace{5\baselineskip}
于是

\nopagebreak[4]
\[
f(x)>f(x_0)-\frac{f(x_0)}2=\frac{f(x_0)}2>0,
\]

所以在邻域 $(x_0-\delta,x_0+\delta)$ 内处处有 $f(x)>0$。

\gkanswer{答案：$|m|<5$ 时有两个交点，$|m|=5$ 时有一个交点，$|m|>5$ 时没有交点；相切时 $m=5$ 或 $m=-5$。}

\end{answer}